\documentclass{article}
\usepackage{amsmath}
\usepackage{amssymb}
\usepackage{booktabs}
\usepackage[a4paper,margin=1in]{geometry}

\title{Model Evaluation Report for T25}
\date{July 18, 2026}
\author{}

\begin{document}
\maketitle

\begin{abstract}
This report summarizes a stratified evaluation (T25) of the promoter-prediction models. Unlike the earlier tasks, T25 introduces no new data: it takes the held-out test predictions already scored in T22 and breaks each model's performance down along three axes --- data source / species, GC-content quartile, and sigma factor --- to locate \emph{where} the models are strong or weak rather than reporting a single aggregate number. Across every stratum the ranking is stable, with ProkBERT strongest, NT\_LoRA close behind, and the CNN a solid but weaker baseline. The models remain strong across species and sigma factors; their one consistent weak spot is the AT-rich GC band, exactly the composition regime in which \emph{Clostridium} promoter motifs and genomic background overlap most.
\end{abstract}

\section{Purpose}
The purpose of this report is to characterize \emph{where} the trained models' held-out performance comes from, by stratifying the T22 test predictions along three axes: data source and species, GC content, and sigma factor. This complements the aggregate test metrics (T22), the stability analysis (T23), and the composition-controlled generalization test (T24) by showing whether performance is uniform across the data or concentrated in --- or degraded on --- particular subpopulations. In particular, it asks whether the AT-composition confound that is central to this project (\emph{Clostridium} is AT-rich, and so is the promoter motif) is visible as a stratum-level weakness in the in-distribution test set.

\section{Confusion matrix and basic quantities}
All the classification metrics considered here are built from the confusion matrix at a given decision threshold. For a binary classification problem, let
\begin{itemize}
\item $TP$: true positives (predicted positive and actually positive),
\item $FP$: false positives (predicted positive but actually negative),
\item $TN$: true negatives (predicted negative and actually negative),
\item $FN$: false negatives (predicted negative but actually positive).
\end{itemize}
From these counts, several basic rates are defined:
\begin{align}
\text{True Positive Rate (TPR)} &= \frac{TP}{TP+FN}, \\
\text{False Positive Rate (FPR)} &= \frac{FP}{FP+TN}, \\
\text{Precision} &= \frac{TP}{TP+FP}, \\
\text{Recall} &= \frac{TP}{TP+FN}.
\end{align}
TPR and recall are the same quantity. Precision measures how many predicted positives are correct, while recall measures how many actual positives are correctly found.

\section{Metric definitions and formulas}

\subsection{AUROC: Area Under the ROC Curve}
The receiver operating characteristic (ROC) curve plots the true positive rate against the false positive rate as the classification threshold is varied. Conceptually, the calculation proceeds as follows:
\begin{enumerate}
\item Sort all examples by their predicted score.
\item Sweep a threshold from high to low over this score range.
\item At each threshold, compute $\mathrm{TPR}=\frac{TP}{TP+FN}$ and $\mathrm{FPR}=\frac{FP}{FP+TN}$.
\item Plot points $(\mathrm{FPR},\mathrm{TPR})$ and connect them to form the ROC curve.
\end{enumerate}
The AUROC is the area under this curve:
\begin{equation}
\mathrm{AUROC} = \int_0^1 \mathrm{TPR}(\mathrm{FPR})\, d(\mathrm{FPR}).
\end{equation}
In practice this integral is approximated numerically on the discrete points obtained from the threshold sweep. An AUROC close to 1 indicates that the model strongly separates positive and negative examples across thresholds, while a value near 0.5 corresponds to random guessing. AUROC is prevalence-independent, so its values are comparable across strata of different positive fractions.

\subsection{Median rank-percentile}
The rank-percentile places each window's score on a common $[0,1]$ scale defined by the whole test set, which makes it directly comparable across strata without any calibration. For predicted scores $\{s_j\}$ over all $N$ test windows, the rank-percentile of window $i$ is
\begin{equation}
\mathrm{pctl}(i) = \frac{\big|\{\, j : s_j \le s_i \,\}\big|}{N}\in(0,1],
\end{equation}
i.e. the fraction of all test windows scoring no higher than window $i$. For each stratum we report the \emph{median} rank-percentile of its positives: a value near $1$ means the stratum's true promoters sit near the top of the genome-wide score ranking, while $0.5$ is the no-skill expectation.

\subsection{Stratified AUROC and omitted metrics}
Within each stratum, AUROC (Eq.~5) is computed between that stratum's positives and the \emph{shared} pool of all test negatives, with a $90\%$ confidence interval obtained by bootstrap resampling of the positives. A complementary AUROC against the SRS-only negatives --- dinucleotide-shuffled real promoters, which preserve base and dinucleotide composition but destroy the $-10/-35$ motif --- is also available as a motif-versus-composition control: a small drop from the all-negative AUROC indicates the stratum is separated by genuine motif structure rather than by base composition. AUPRC, F1, and recall at a fixed $0.5$ threshold are deliberately omitted here, because the AUPRC baseline equals a stratum's positive fraction (and so varies between strata, making cross-stratum comparison misleading), and a $0.5$ threshold is not meaningful on an uncalibrated score trained under class balancing.

\section{T25: Stratified evaluation of the held-out test set}

\subsection{Setup}
Task T25 does not train or score any new model; it re-uses the held-out T15 \emph{test} predictions already produced in T22 (row-aligned with the T15 test metadata) and stratifies them. The test set contains 3,323 windows --- 969 verified positives and 2,354 negatives (a shared pool) --- all $81$\,bp. Each model's positives are grouped three ways: by data source / species, by GC-content quartile, and by sigma factor. For every stratum we report the median rank-percentile of its positives and the AUROC of those positives against the shared negative pool, with a $90\%$ bootstrap confidence interval.

The sigma-factor axis is the weakest of the three and is reported with caveats. Only 339 of the 969 positives carry a sigma label, and every label is cross-species: RegulonDB (\emph{E.\ coli}) supplies real per-promoter sigma assignments, while the DBTBS (\emph{B.\ subtilis}) set is labelled with a single blanket ``SigmaA''. There are no in-genus sigma labels, and the composition-relevant class $\sigma^{54}$ has only $n=6$ positives in this split. Sigma is therefore a weak, species-confounded proxy, and sigma-versus-sigma comparisons are only valid within RegulonDB; the GC-content quartile is the dense, in-genus-inclusive version of the same composition question.

\subsection{Results}
Table~1 gives the overall (unstratified) held-out performance, which matches the aggregate T22 test AUROCs as expected --- T25 re-scores the same predictions.

\begin{table}[h]
\centering
\caption{T25 overall: held-out test performance (969 verified positives vs.\ 2{,}354 shared negatives).}
\begin{tabular}{lcc}
\toprule
Model & Median rank-percentile & AUROC [90\% CI] \\
\midrule
ProkBERT & 0.848 & 0.944 [0.94--0.95] \\
NT\_LoRA  & 0.841 & 0.935 [0.93--0.94] \\
CNN      & 0.834 & 0.909 [0.90--0.92] \\
\bottomrule
\end{tabular}
\end{table}

All three models place a typical real promoter above roughly $84$--$85\%$ of all test windows, and the overall ordering (ProkBERT $>$ NT\_LoRA $>$ CNN) is preserved in every stratum below.

\begin{table}[h]
\centering
\caption{T25 by data source / species (AUROC vs.\ the shared negative pool).}
\begin{tabular}{lccc}
\toprule
Stratum ($n_{\text{pos}}$) & ProkBERT & NT\_LoRA & CNN \\
\midrule
Hocq 2022 in-genus TSS (573)   & 0.946 & 0.923 & 0.898 \\
RegulonDB \emph{E.\ coli} (358) & 0.942 & 0.956 & 0.932 \\
DBTBS \emph{B.\ subtilis} (38) & 0.922 & 0.918 & 0.855 \\
\bottomrule
\end{tabular}
\end{table}

Every source is scored well (AUROC $\ge 0.90$ for the two transformers). Notably, ProkBERT's \emph{best} source is the in-genus Hocq 2022 TSS (0.946), so its in-distribution performance is not carried by the cross-species \emph{E.\ coli} data. NT\_LoRA shows the opposite pattern --- strongest on RegulonDB \emph{E.\ coli} (0.956) and weaker on the in-genus set (0.923) --- consistent with it leaning more on the cross-species signal.

\begin{table}[h]
\centering
\caption{T25 by GC-content quartile (AUROC vs.\ the shared negative pool). \emph{Clostridium} promoters are AT-rich (GC $\approx 28$--$32\%$).}
\begin{tabular}{lccc}
\toprule
GC quartile ($n_{\text{pos}}$) & ProkBERT & NT\_LoRA & CNN \\
\midrule
7--21\% (273)  & 0.954 & 0.940 & 0.920 \\
21--26\% (215) & 0.954 & 0.933 & 0.905 \\
26--37\% (240) & 0.916 & 0.894 & 0.859 \\
37--65\% (241) & 0.951 & 0.973 & 0.948 \\
\bottomrule
\end{tabular}
\end{table}

This is the diagnostic axis. Every model dips in the $26$--$37\%$ GC band --- the AT-rich range where promoter composition overlaps the genomic background most --- and nowhere else. ProkBERT falls from $\sim$0.95 to 0.916, the CNN from $\sim$0.95 to 0.859, and NT\_LoRA dips hardest, from 0.973 in the highest-GC band to 0.894. The median rank-percentile shows the same dip (ProkBERT 0.80 in this band versus $\sim$0.87 in the others).

\begin{table}[h]
\centering
\caption{T25 by sigma factor (AUROC vs.\ the shared negative pool; cross-species labels only --- see Setup).}
\begin{tabular}{lccc}
\toprule
Sigma factor ($n_{\text{pos}}$) & ProkBERT & NT\_LoRA & CNN \\
\midrule
$\sigma^{70}$ (174) & 0.942 & 0.956 & 0.932 \\
$\sigma^{24}$ (49)  & 0.922 & 0.939 & 0.910 \\
\emph{B.\ subtilis} / SigmaA (38) & 0.922 & 0.918 & 0.855 \\
$\sigma^{32}$ (29)  & 0.947 & 0.976 & 0.956 \\
$\sigma^{38}$ (28)  & 0.924 & 0.965 & 0.939 \\
$\sigma^{28}$ (15)  & 0.939 & 0.948 & 0.958 \\
$\sigma^{54}$ (6)$^{\dagger}$ & 0.986 & 0.946 & 0.942 \\
unlabeled, mostly in-genus T7b (630) & 0.947 & 0.926 & 0.900 \\
\bottomrule
\end{tabular}

{\footnotesize $^{\dagger}$ $n<10$ positives; shown for completeness, not a reliable per-stratum estimate.}
\end{table}

Performance is essentially uniform across sigma classes (AUROC $\approx 0.92$--$0.98$ for all adequately-powered strata), with no class systematically missed. Given the caveats above --- cross-species-only labels, species confounding, and the tiny $\sigma^{54}$ stratum --- the honest reading is that no reliable sigma-specific effect is detectable, and the GC-content axis is the meaningful version of this question.

\subsection{Interpretation}
Three findings stand out.

First, the single consistent weak spot is the $26$--$37\%$ GC band. It is the only stratum on any axis where all three models drop together (ProkBERT 0.916, NT\_LoRA 0.894, CNN 0.859), and it corresponds exactly to the AT-rich composition where \emph{Clostridium} promoter motifs and genomic background overlap. This is the AT-composition confound --- the core risk of the project --- appearing as a real but \emph{modest} in-distribution weakness. NT\_LoRA dips hardest here, marking it as the most composition-reliant of the three; this is the same weakness that becomes a full collapse to chance for NT\_LoRA on the out-of-distribution, composition-matched test (T24).

Second, in-genus data is not a weak stratum. ProkBERT's strongest source is the in-genus Hocq 2022 TSS, so in-distribution the models generalize to the target genus without penalty; the difficulty seen in T24 is specific to composition-matched negatives and the out-of-genus subset, not to in-genus positives as such.

Third, the model ordering is robust. ProkBERT $\ge$ NT\_LoRA $\ge$ CNN holds on all three axes with tight confidence intervals, so the aggregate ranking from T22/T23 is not an artifact of one particular split or subpopulation.

Taken together, T25 localizes the models' only systematic weakness to the AT-rich composition band and confirms strength everywhere else. This is precisely the weakness that the downstream calibration steps are designed to correct: an AT-content--aware calibration (T26) and an empirical FDR built from the SRS composition-matched null (T27).

\section{Overall conclusion}
T25 stratifies the held-out test predictions rather than adding new data, and it gives a consistent picture: all three models are strong and well-ranked across data sources, species, and sigma factors, and their performance is not carried by any single subpopulation. The one systematic weakness is the AT-rich GC band, where every model dips and NT\_LoRA dips most --- the in-distribution shadow of the composition confound that turns into NT\_LoRA's chance-level result on the composition-matched T24 test. ProkBERT is the strongest and most composition-robust model, NT\_LoRA follows closely in-distribution but is the most composition-sensitive, and the CNN is a consistent but weaker baseline. The practical implication is unchanged from T24: the priority for improving genome-wide precision is reducing composition-reliance --- through AT-aware calibration (T26), an SRS-null empirical FDR (T27), and iterative hard-negative mining --- concentrated on the AT-rich band that T25 identifies.

\end{document}
